% GENERATED FILE. Edit canonical lessons, then run npm run generate:books.
% canonical-source-hash: fnv1a64:4c9460c05c1082ce
% canonical-lessons: TE-C484-ammakam, TE-C484-viniyogadarudu, TE-C484-nanyata, TE-C484-rayiti, TE-C484-beramadu

\chapter{Buying and Selling}
\label{ch:te-buying-and-selling}
\hlchaptermodality{484}

\begin{chapteropening}
\textbf{By the end of this chapter:} \emph{I can say a sale, a consumer, quality, a concession and to haggle.}

Complete the last lesson of chapter 484: I can say a sale, a consumer, quality, a concession and to haggle.
\end{chapteropening}

\section[\texorpdfstring{\te{అమ్మకం} (ammakaṁ)}{అమ్మకం (ammakaṁ)}]{\te{అమ్మకం} (ammakaṁ) — a sale}

\label{lesson:TE-C484-ammakam}



\begin{quote}
Before the new one: say the Telugu for to go out, then the Telugu for to put forth leaves.
\end{quote}

\subsection*{You'll want to know: \te{అమ్మకం}}
\textbf{\te{అమ్మకం}} — \emph{ammakaṁ} — \textquotedblleft{}a sale\textquotedblright{}.

\begin{quote}
\te{పండుగ} \te{ముందు} \te{బట్టల} \te{అమ్మకాలు} \te{పెరిగాయి}.
\end{quote}

\begin{quote}
\emph{paṇḍuga mundu baṭṭala ammakālu perigāyi.}
\end{quote}

\begin{quote}
Before the festival the sales of clothes went up.
\end{quote}

Plural: \te{అమ్మకాలు}, \emph{ammakālu}.

The noun of \te{అమ్ము}, to sell. Plural \te{అమ్మకాలు}, sales.

\begin{cousinweb}
\textbf{\te{అమ్మకం}} (\emph{ammakaṁ}) is native Dravidian: DEDR 186 lists Telugu \emph{ammakamu}, \textquotedblleft{}sale\textquotedblright{}, beside \emph{ammu}, \textquotedblleft{}to sell\textquotedblright{}.
\end{cousinweb}

\begin{grammarlens}[title={What you've built}]
One more word for getting around.
\end{grammarlens}

\subsection*{Your turn}
\begin{itemize}
\raggedright
  \item \emph{Say it:} \emph{ammakaṁ}
  \item \emph{Say it:} \emph{ammakaṁ}, once more
  \item \emph{Say it:} \emph{paṇḍuga mundu baṭṭala ammakālu perigāyi.} — Before the festival the sales of clothes went up.
  \item \emph{Recall:} say the Telugu for to go out, then the Telugu for to put forth leaves, then say \emph{ammakaṁ} again
\end{itemize}

\begin{tcolorbox}[breakable,colback=teal!4,colframe=teal!35!black,title={Before you move on}]
What does \te{అమ్మకం} mean? (\textquotedblleft{}A sale\textquotedblright{}.) Say it once more.
\end{tcolorbox}
\section[\texorpdfstring{\te{వినియోగదారుడు} (viniyōgadāruḍu)}{వినియోగదారుడు (viniyōgadāruḍu)}]{\te{వినియోగదారుడు} (viniyōgadāruḍu) — a consumer; a customer}

\label{lesson:TE-C484-viniyogadarudu}



\begin{quote}
Before the new one: say the Telugu for to put forth leaves, then the Telugu for a sale.
\end{quote}

\subsection*{You'll want to know: \te{వినియోగదారుడు}}
\textbf{\te{వినియోగదారుడు}} — \emph{viniyōgadāruḍu} — \textquotedblleft{}a consumer; a customer\textquotedblright{}.

\begin{quote}
\te{వినియోగదారుడికి} \te{రసీదు} \te{అడిగే} \te{హక్కు} \te{ఉంది}.
\end{quote}

\begin{quote}
\emph{viniyōgadāruḍiki rasīdu aḍigē hakku undi.}
\end{quote}

\begin{quote}
A consumer has the right to ask for a receipt.
\end{quote}

Plural: \te{వినియోగదారులు}, \emph{viniyōgadārulu}.

\te{వినియోగం}, use, + -\te{దారుడు}, one who has: one who uses goods. Plural \te{వినియోగదారులు}.

\begin{cousinweb}
\textbf{\te{వినియోగదారుడు}} (\emph{viniyōgadāruḍu}) is built from Sanskrit \emph{viniyoga}, \textquotedblleft{}employment, use\textquotedblright{} (Monier-Williams), and -\te{దారుడు}, \textquotedblleft{}one who holds\textquotedblright{}, from Persian \emph{-dār}, \textquotedblleft{}holder\textquotedblright{}.
\end{cousinweb}

\begin{grammarlens}[title={What you've built}]
One more word for getting around.
\end{grammarlens}

\subsection*{Your turn}
\begin{itemize}
\raggedright
  \item \emph{Say it:} \emph{viniyōgadāruḍu}
  \item \emph{Say it:} \emph{viniyōgadāruḍu}, once more
  \item \emph{Say it:} \emph{viniyōgadāruḍiki rasīdu aḍigē hakku undi.} — A consumer has the right to ask for a receipt.
  \item \emph{Recall:} say the Telugu for to put forth leaves, then the Telugu for a sale, then say \emph{viniyōgadāruḍu} again
\end{itemize}

\begin{tcolorbox}[breakable,colback=teal!4,colframe=teal!35!black,title={Before you move on}]
What does \te{వినియోగదారుడు} mean? (\textquotedblleft{}A consumer; a customer\textquotedblright{}.) Say it once more.
\end{tcolorbox}
\section[\texorpdfstring{\te{నాణ్యత} (nāṇyata)}{నాణ్యత (nāṇyata)}]{\te{నాణ్యత} (nāṇyata) — quality}

\label{lesson:TE-C484-nanyata}



\begin{quote}
Before the new one: say the Telugu for a sale, then the Telugu for a consumer.
\end{quote}

\subsection*{You'll want to know: \te{నాణ్యత}}
\textbf{\te{నాణ్యత}} — \emph{nāṇyata} — \textquotedblleft{}quality\textquotedblright{}.

\begin{quote}
\te{ఆ} \te{దుకాణంలో} \te{బట్టల} \te{నాణ్యత} \te{బాగుంటుంది}.
\end{quote}

\begin{quote}
\emph{ā dukāṇaṁlō baṭṭala nāṇyata bāguṇṭundi.}
\end{quote}

\begin{quote}
The quality of the clothes in that shop is good.
\end{quote}

\te{నాణ్యమైన}, of good quality: \te{నాణ్యమైన} \te{సరుకు}, good goods.

\begin{cousinweb}
\textbf{\te{నాణ్యత}} (\emph{nāṇyata}) is a Telugu word for quality whose origin is not settled.
\end{cousinweb}

\begin{grammarlens}[title={What you've built}]
One more word for getting around.
\end{grammarlens}

\subsection*{Your turn}
\begin{itemize}
\raggedright
  \item \emph{Say it:} \emph{nāṇyata}
  \item \emph{Say it:} \emph{nāṇyata}, once more
  \item \emph{Say it:} \emph{ā dukāṇaṁlō baṭṭala nāṇyata bāguṇṭundi.} — The quality of the clothes in that shop is good.
  \item \emph{Recall:} say the Telugu for a sale, then the Telugu for a consumer, then say \emph{nāṇyata} again
\end{itemize}

\begin{tcolorbox}[breakable,colback=teal!4,colframe=teal!35!black,title={Before you move on}]
What does \te{నాణ్యత} mean? (\textquotedblleft{}Quality\textquotedblright{}.) Say it once more.
\end{tcolorbox}
\section[\texorpdfstring{\te{రాయితీ} (rāyitī)}{రాయితీ (rāyitī)}]{\te{రాయితీ} (rāyitī) — a concession; a rebate}

\label{lesson:TE-C484-rayiti}



\begin{quote}
Before the new one: say the Telugu for a consumer, then the Telugu for quality.
\end{quote}

\subsection*{You'll want to know: \te{రాయితీ}}
\textbf{\te{రాయితీ}} — \emph{rāyitī} — \textquotedblleft{}a concession; a rebate\textquotedblright{}.

\begin{quote}
\te{విద్యార్థులకు} \te{టిక్కెట్టు} \te{మీద} \te{రాయితీ} \te{ఉంది}.
\end{quote}

\begin{quote}
\emph{vidyārthulaku ṭikkeṭṭu mīda rāyitī undi.}
\end{quote}

\begin{quote}
Students get a concession on the ticket.
\end{quote}

A reduction given to a group: students, old people, farmers.

\begin{cousinweb}
\textbf{\te{రాయితీ}} (\emph{rāyitī}) is usually traced to Arabic \emph{riʿāya}, \textquotedblleft{}regard, consideration\textquotedblright{}, which came through Persian and Hindustani.
\end{cousinweb}

\begin{grammarlens}[title={What you've built}]
One more word for getting around.
\end{grammarlens}

\subsection*{Your turn}
\begin{itemize}
\raggedright
  \item \emph{Say it:} \emph{rāyitī}
  \item \emph{Say it:} \emph{rāyitī}, once more
  \item \emph{Say it:} \emph{vidyārthulaku ṭikkeṭṭu mīda rāyitī undi.} — Students get a concession on the ticket.
  \item \emph{Recall:} say the Telugu for a consumer, then the Telugu for quality, then say \emph{rāyitī} again
\end{itemize}

\begin{tcolorbox}[breakable,colback=teal!4,colframe=teal!35!black,title={Before you move on}]
What does \te{రాయితీ} mean? (\textquotedblleft{}A concession; a rebate\textquotedblright{}.) Say it once more.
\end{tcolorbox}
\section[\texorpdfstring{\te{బేరమాడు} (bēramāḍu)}{బేరమాడు (bēramāḍu)}]{\te{బేరమాడు} (bēramāḍu) — to haggle}

\label{lesson:TE-C484-beramadu}



\begin{quote}
Before the new one: say the Telugu for quality, then the Telugu for a concession.
\end{quote}

\subsection*{You'll want to know: \te{బేరమాడు}}
\textbf{\te{బేరమాడు}} — \emph{bēramāḍu} — \textquotedblleft{}to haggle\textquotedblright{}.

\begin{quote}
\te{అమ్మ} \te{కూరగాయల} \te{ధర} \te{కోసం} \te{బేరమాడింది}.
\end{quote}

\begin{quote}
\emph{amma kūragāyala dhara kōsaṁ bēramāḍindi.}
\end{quote}

\begin{quote}
Mother haggled over the price of the vegetables.
\end{quote}

Past (I …): \te{బేరమాడాను}, \emph{bēramāḍānu}.

\te{బేరం}, a bargain, + \te{ఆడు}, to play: to play at bargaining.

\begin{cousinweb}
\textbf{\te{బేరమాడు}} (\emph{bēramāḍu}) is built on \te{బేరం}, \textquotedblleft{}a bargain\textquotedblright{}, whose origin is not settled, with the native \te{ఆడు}, \textquotedblleft{}to play, to do\textquotedblright{}.
\end{cousinweb}

\begin{grammarlens}[title={What you've built}]
One more word for getting around.
\end{grammarlens}

\subsection*{Your turn}
\begin{itemize}
\raggedright
  \item \emph{Say it:} \emph{bēramāḍu}
  \item \emph{Say it:} \emph{bēramāḍu}, once more
  \item \emph{Say it:} \emph{amma kūragāyala dhara kōsaṁ bēramāḍindi.} — Mother haggled over the price of the vegetables.
  \item \emph{Recall:} say the Telugu for quality, then the Telugu for a concession, then say \emph{bēramāḍu} again
\end{itemize}

\begin{tcolorbox}[breakable,colback=teal!4,colframe=teal!35!black,title={Before you move on}]
What does \te{బేరమాడు} mean? (\textquotedblleft{}To haggle\textquotedblright{}.) Say it once more.
\end{tcolorbox}
